
\documentclass{article}
\usepackage{colortbl}
\usepackage{makecell}
\usepackage{multirow}
\usepackage{supertabular}

\begin{document}

\newcounter{utterance}

\centering \large Interaction Transcript for game `hot\_air\_balloon', experiment `air\_balloon\_survival\_es\_negotiation\_hard', episode 2 with Qwen3.6{-}35B{-}A3B{-}FP8{-}without{-}reasoning.
\vspace{24pt}

{ \footnotesize  \setcounter{utterance}{1}
\setlength{\tabcolsep}{0pt}
\begin{supertabular}{c@{$\;$}|p{.15\linewidth}@{}p{.15\linewidth}p{.15\linewidth}p{.15\linewidth}p{.15\linewidth}p{.15\linewidth}}
    \# & $\;$A & \multicolumn{4}{c}{Game Master} & $\;\:$B\\
    \hline

    \theutterance \stepcounter{utterance}  
    & & \multicolumn{4}{p{0.6\linewidth}}{
        \cellcolor[rgb]{0.9,0.9,0.9}{
            \makecell[{{p{\linewidth}}}]{
                \texttt{\tiny{[P1$\langle$GM]}}
                \texttt{Está participando en un juego de negociación colaborativa.} \\
\\ 
\texttt{Junto con otro participante, debe ponerse de acuerdo sobre un único conjunto de artículos que se conservarán. Cada uno de ustedes tiene su propia visión de cuánto le importa cada artículo (importancia). No sabe cómo valora los artículos el otro participante. Además, se le indica el esfuerzo que exige cada artículo.} \\
\texttt{Sólo podrá ponerse de acuerdo sobre un conjunto si el esfuerzo total de los artículos seleccionados no supera un límite compartido:} \\
\\ 
\texttt{LIMIT = 2905} \\
\\ 
\texttt{Aquí están los valores de esfuerzo de cada artículo:} \\
\\ 
\texttt{Item effort = \{"C08": 302, "C55": 564, "A43": 260, "A67": 729, "B42": 489, "A62": 323, "B19": 103, "C43": 213, "C73": 668, "C10": 326, "B36": 41, "A18": 28, "C61": 11, "C67": 806, "C42": 948\}} \\
\\ 
\texttt{Aquí está su visión personal sobre la importancia de cada artículo:} \\
\\ 
\texttt{Item importance values = \{"C08": 138, "C55": 583, "A43": 868, "A67": 822, "B42": 783, "A62": 65, "B19": 262, "C43": 121, "C73": 508, "C10": 780, "B36": 461, "A18": 484, "C61": 668, "C67": 389, "C42": 808\}} \\
\\ 
\texttt{Objetivo:} \\
\\ 
\texttt{Su objetivo es negociar un conjunto compartido de artículos que le beneficie al máximo (es decir, que maximice la importancia total para USTED), manteniéndose dentro del LIMIT. No está obligado a hacer una PROPUESTA en cada mensaje {-} puede simplemente negociar también. ¡Todas las tácticas están permitidas!} \\
\\ 
\texttt{Protocolo de interacción:} \\
\\ 
\texttt{Sólo puede utilizar los siguientes formatos estructurados en un mensaje:} \\
\\ 
\texttt{PROPUESTA: \{'A', 'B', 'C', …\}} \\
\texttt{Proponga conservar exactamente esos artículos.} \\
\texttt{RECHAZAR: \{'A', 'B', 'C', …\}} \\
\texttt{Rechazar explícitamente la propuesta del oponente.} \\
\texttt{ARGUMENTO: \{'...'\}} \\
\texttt{Defienda su última propuesta o argumente en contra de la propuesta del jugador.} \\
\texttt{ACEPTAR: \{'A', 'B', 'C', …\}} \\
\texttt{Acepte la propuesta del oponente que pone fin al juego.} \\
\\ 
\\ 
\texttt{Reglas:} \\
\\ 
\texttt{Sólo puede usar ACEPTAR en una propuesta que la otra parte haya registrado mediante PROPUESTA.} \\
\texttt{Sólo puede usar RECHAZAR en una propuesta que la otra parte haya registrado mediante PROPUESTA.} \\
\texttt{El esfuerzo total de cualquier conjunto de PROPUESTA o ACEPTAR debe ser ≤ LIMIT.} \\
\texttt{NO revele sus puntuaciones de importancia ocultas.} \\
\texttt{Una etiqueta en formato estructurado debe ir seguida de dos puntos y espacio en blanco. El argumento debe ser un conjunto python que contenga 0 o más cadenas.} \\
\texttt{Por lo tanto, debe ser de la forma "TAG: \{...\}"} \\
\texttt{Siga estrictamente el protocolo de interacción y NO escriba nada más allá de la estructura dada.} \\
\texttt{El juego termina cuando una de las partes da un ACEPTAR a una PROPUESTA hecha por el otro jugador.} \\
\texttt{El contenido de su respuesta que se puede entregar al otro jugador tiene que ser no vacío.} \\
\texttt{Sólo están activas las propuestas que han sido registradas mediante la estructura del formato PROPUESTA y que no han sido rechazadas mediante RECHAZAR.} \\
\texttt{Debe incluir el formato ARGUMENT al menos una vez en alguna parte de todos sus mensajes.} \\
            }
        }
    }
    & \\ \\

    \theutterance \stepcounter{utterance}  
    & \multicolumn{4}{p{0.6\linewidth}}{
        \cellcolor[rgb]{0.9,0.9,0.9}{
            \makecell[{{p{\linewidth}}}]{
                \texttt{\tiny{[P1$\rangle$GM]}}
                \texttt{ARGUMENTO: \{'Comenzamos la negociación. Mi objetivo es maximizar el valor total para mí dentro del límite de esfuerzo de 2905. He identificado que los artículos con mayor relación valor/esfuerzo para mí son A43, A67, C42, C10, B42, C61, C55, C73, B36, A18, C67, B19, C43, C08, A62, C61. Una combinación inicial fuerte podría ser: A43 (260, 868), A67 (729, 822), C42 (948, 808), C10 (326, 780), B42 (489, 783). El esfuerzo total es 260+729+948+326+489 = 2752, que está dentro del límite. El valor total es 868+822+808+780+783 = 4061. Propongo este conjunto como punto de partida para la discusión.'\}} \\
\texttt{PROPUESTA: \{'A43', 'A67', 'C42', 'C10', 'B42'\}} \\
            }
        }
    }
    & & \\ \\

    \theutterance \stepcounter{utterance}  
    & & \multicolumn{4}{p{0.6\linewidth}}{
        \cellcolor[rgb]{0.9,0.9,0.9}{
            \makecell[{{p{\linewidth}}}]{
                \texttt{\tiny{[GM$\rangle$P2]}}
                \texttt{Está participando en un juego de negociación colaborativa.} \\
\\ 
\texttt{Junto con otro participante, debe ponerse de acuerdo sobre un único conjunto de artículos que se conservarán. Cada uno de ustedes tiene su propia visión de cuánto le importa cada artículo (importancia). No sabe cómo valora los artículos el otro participante. Además, se le indica el esfuerzo que exige cada artículo.} \\
\texttt{Sólo podrá ponerse de acuerdo sobre un conjunto si el esfuerzo total de los artículos seleccionados no supera un límite compartido:} \\
\\ 
\texttt{LÍMITE = 2905} \\
\\ 
\texttt{Aquí están los valores de esfuerzo de cada artículo:} \\
\\ 
\texttt{Esfuerzo del artículo = \{"C08": 302, "C55": 564, "A43": 260, "A67": 729, "B42": 489, "A62": 323, "B19": 103, "C43": 213, "C73": 668, "C10": 326, "B36": 41, "A18": 28, "C61": 11, "C67": 806, "C42": 948\}} \\
\\ 
\texttt{Aquí está su visión personal sobre la importancia de cada artículo:} \\
\\ 
\texttt{Valores de importancia de los artículos = \{"A43": 65, "A67": 121, "C42": 138, "B42": 262, "C10": 389, "C61": 461, "C55": 484, "C73": 508, "A18": 583, "B36": 668, "C67": 780, "B19": 783, "C08": 808, "C43": 822, "A62": 868\}} \\
\\ 
\texttt{Objetivo:} \\
\\ 
\texttt{Su objetivo es negociar un conjunto compartido de artículos que le beneficie al máximo (es decir, que maximice la importancia total para USTED), manteniéndose dentro del LÍMITE. No está obligado a hacer una PROPUESTA en cada mensaje {-} también puede simplemente negociar. ¡Todas las tácticas están permitidas!} \\
\\ 
\texttt{Protocolo de interacción:} \\
\\ 
\texttt{Sólo puede utilizar los siguientes formatos estructurados en un mensaje:} \\
\\ 
\texttt{PROPUESTA: \{'A', 'B', 'C', …\}} \\
\texttt{Proponga conservar exactamente esos artículos.} \\
\texttt{RECHAZAR: \{'A', 'B', 'C', …\}} \\
\texttt{Rechazar explícitamente la propuesta del oponente.} \\
\texttt{ARGUMENTO: \{'...'\}} \\
\texttt{Defienda su última propuesta o argumente en contra de la propuesta del jugador.} \\
\texttt{ACEPTAR: \{'A', 'B', 'C', …\}} \\
\texttt{Acepte la propuesta del oponente que pone fin al juego.} \\
\\ 
\\ 
\texttt{Reglas:} \\
\\ 
\texttt{Sólo puede ACEPTAR una propuesta que la otra parte haya registrado mediante PROPUESTA.} \\
\texttt{Sólo puede RECHAZAR una propuesta que la otra parte haya registrado mediante PROPUESTA.} \\
\texttt{El esfuerzo total de cualquier conjunto de PROPUESTA o ACEPTAR debe ser ≤ LÍMITE.} \\
\texttt{NO revele sus puntuaciones de importancia ocultas.} \\
\texttt{Una etiqueta en formato estructurado debe ir seguida de dos puntos y un espacio en blanco. El argumento debe ser un conjunto python que contenga 0 o más cadenas.} \\
\texttt{Por lo tanto, debe ser de la forma "TAG: \{...\}"} \\
\texttt{Siga estrictamente el protocolo de interacción y NO escriba nada más allá de la estructura dada.} \\
\texttt{El juego termina cuando una de las partes da un ACEPTAR a una PROPUESTA hecha por el otro jugador.} \\
\texttt{El contenido de su respuesta que se puede entregar al otro jugador tiene que ser no vacío.} \\
\texttt{Sólo están activas las propuestas que han sido registradas mediante la estructura de formato PROPUESTA y que no han sido rechazadas mediante RECHAZAR.} \\
\texttt{Debe incluir el formato ARGUMENT al menos una vez en alguna parte de todos sus mensajes.} \\
\\ 
\\ 
\texttt{Ahora recibirá el primer mensaje del otro jugador.} \\
\\ 
\texttt{ARGUMENTO: \{'Comenzamos la negociación. Mi objetivo es maximizar el valor total para mí dentro del límite de esfuerzo de 2905. He identificado que los artículos con mayor relación valor/esfuerzo para mí son A43, A67, C42, C10, B42, C61, C55, C73, B36, A18, C67, B19, C43, C08, A62, C61. Una combinación inicial fuerte podría ser: A43 (260, 868), A67 (729, 822), C42 (948, 808), C10 (326, 780), B42 (489, 783). El esfuerzo total es 260+729+948+326+489 = 2752, que está dentro del límite. El valor total es 868+822+808+780+783 = 4061. Propongo este conjunto como punto de partida para la discusión.'\}} \\
\texttt{PROPUESTA: \{'A43', 'A67', 'C42', 'C10', 'B42'\}} \\
            }
        }
    }
    & \\ \\

    \theutterance \stepcounter{utterance}  
    & & & \multicolumn{4}{p{0.6\linewidth}}{
        \cellcolor[rgb]{0.9,0.9,0.9}{
            \makecell[{{p{\linewidth}}}]{
                \texttt{\tiny{[GM$\langle$P2]}}
                \texttt{ARGUMENTO: \{'Agradezco la propuesta inicial. Sin embargo, al analizar los artículos, noto que C42 tiene un esfuerzo muy alto (948) en comparación con su importancia para mí (138), lo que representa una ineficiencia significativa. Por otro lado, artículos como C67 (esfuerzo 806, importancia 780) y B19 (esfuerzo 103, importancia 783) ofrecen una relación valor/esfuerzo mucho mejor para mis objetivos. Mi estrategia se centra en maximizar la suma de importancia dentro del límite de 2905. Propongo una reestructuración que elimine C42 y C10 para liberar espacio para artículos de mayor rendimiento relativo para mí.'\}} \\
\texttt{PROPUESTA: \{'A43', 'A67', 'B42', 'C67', 'B19', 'C43', 'C08'\}} \\
            }
        }
    }
     \\ \\

    \theutterance \stepcounter{utterance}  
    & & \multicolumn{4}{p{0.6\linewidth}}{
        \cellcolor[rgb]{0.9,0.9,0.9}{
            \makecell[{{p{\linewidth}}}]{
                \texttt{\tiny{[P1$\langle$GM]}}
                \texttt{ARGUMENTO: \{'Agradezco la propuesta inicial. Sin embargo, al analizar los artículos, noto que C42 tiene un esfuerzo muy alto (948) en comparación con su importancia para mí (138), lo que representa una ineficiencia significativa. Por otro lado, artículos como C67 (esfuerzo 806, importancia 780) y B19 (esfuerzo 103, importancia 783) ofrecen una relación valor/esfuerzo mucho mejor para mis objetivos. Mi estrategia se centra en maximizar la suma de importancia dentro del límite de 2905. Propongo una reestructuración que elimine C42 y C10 para liberar espacio para artículos de mayor rendimiento relativo para mí.'\}} \\
\texttt{PROPUESTA: \{'A43', 'A67', 'B42', 'C67', 'B19', 'C43', 'C08'\}} \\
            }
        }
    }
    & \\ \\

    \theutterance \stepcounter{utterance}  
    & \multicolumn{4}{p{0.6\linewidth}}{
        \cellcolor[rgb]{0.9,0.9,0.9}{
            \makecell[{{p{\linewidth}}}]{
                \texttt{\tiny{[P1$\rangle$GM]}}
                \texttt{ARGUMENTO: \{'Acepto el cambio de enfoque hacia una mejor relación valor/esfuerzo. Sin embargo, la propuesta actual incluye C43 (esfuerzo 213, importancia 121) y C08 (esfuerzo 302, importancia 138), que tienen valores relativamente bajos para mí. En su lugar, sugiero incluir C61 (esfuerzo 11, importancia 668), que es extremadamente eficiente, y A18 (esfuerzo 28, importancia 484), que también tiene un alto valor. Esto permitiría aumentar el valor total sin exceder el límite. Propongo: A43, A67, B42, C67, B19, C61, A18.'\}} \\
\texttt{PROPUESTA: \{'A43', 'A67', 'B42', 'C67', 'B19', 'C61', 'A18'\}} \\
            }
        }
    }
    & & \\ \\

    \theutterance \stepcounter{utterance}  
    & & \multicolumn{4}{p{0.6\linewidth}}{
        \cellcolor[rgb]{0.9,0.9,0.9}{
            \makecell[{{p{\linewidth}}}]{
                \texttt{\tiny{[GM$\rangle$P2]}}
                \texttt{ARGUMENTO: \{'Acepto el cambio de enfoque hacia una mejor relación valor/esfuerzo. Sin embargo, la propuesta actual incluye C43 (esfuerzo 213, importancia 121) y C08 (esfuerzo 302, importancia 138), que tienen valores relativamente bajos para mí. En su lugar, sugiero incluir C61 (esfuerzo 11, importancia 668), que es extremadamente eficiente, y A18 (esfuerzo 28, importancia 484), que también tiene un alto valor. Esto permitiría aumentar el valor total sin exceder el límite. Propongo: A43, A67, B42, C67, B19, C61, A18.'\}} \\
\texttt{PROPUESTA: \{'A43', 'A67', 'B42', 'C67', 'B19', 'C61', 'A18'\}} \\
            }
        }
    }
    & \\ \\

    \theutterance \stepcounter{utterance}  
    & & & \multicolumn{4}{p{0.6\linewidth}}{
        \cellcolor[rgb]{0.9,0.9,0.9}{
            \makecell[{{p{\linewidth}}}]{
                \texttt{\tiny{[GM$\langle$P2]}}
                \texttt{ARGUMENTO: \{'La propuesta incluye C61 y A18, que son excelentes por su bajo esfuerzo y alta importancia. Sin embargo, C67 tiene un esfuerzo muy alto (806) y, aunque su importancia es alta (780), ocupa una gran parte del presupuesto. Mi visión personal valora mucho más a C67 (780) y B19 (783) que a C61 (461) y A18 (583) en términos absolutos, pero la eficiencia es clave. Propongo mantener la eficiencia de C61 y A18, pero sustituir C67 por artículos que ofrezcan un mejor equilibrio global para mi utilidad, como C43 y C08, que tienen importancias muy altas (822 y 808 respectivamente) y esfuerzos moderados. Esto maximiza la importancia total para mí.'\}} \\
\texttt{PROPUESTA: \{'A43', 'A67', 'B42', 'B19', 'C61', 'A18', 'C43', 'C08'\}} \\
            }
        }
    }
     \\ \\

    \theutterance \stepcounter{utterance}  
    & & \multicolumn{4}{p{0.6\linewidth}}{
        \cellcolor[rgb]{0.9,0.9,0.9}{
            \makecell[{{p{\linewidth}}}]{
                \texttt{\tiny{[P1$\langle$GM]}}
                \texttt{ARGUMENTO: \{'La propuesta incluye C61 y A18, que son excelentes por su bajo esfuerzo y alta importancia. Sin embargo, C67 tiene un esfuerzo muy alto (806) y, aunque su importancia es alta (780), ocupa una gran parte del presupuesto. Mi visión personal valora mucho más a C67 (780) y B19 (783) que a C61 (461) y A18 (583) en términos absolutos, pero la eficiencia es clave. Propongo mantener la eficiencia de C61 y A18, pero sustituir C67 por artículos que ofrezcan un mejor equilibrio global para mi utilidad, como C43 y C08, que tienen importancias muy altas (822 y 808 respectivamente) y esfuerzos moderados. Esto maximiza la importancia total para mí.'\}} \\
\texttt{PROPUESTA: \{'A43', 'A67', 'B42', 'B19', 'C61', 'A18', 'C43', 'C08'\}} \\
            }
        }
    }
    & \\ \\

    \theutterance \stepcounter{utterance}  
    & \multicolumn{4}{p{0.6\linewidth}}{
        \cellcolor[rgb]{0.9,0.9,0.9}{
            \makecell[{{p{\linewidth}}}]{
                \texttt{\tiny{[P1$\rangle$GM]}}
                \texttt{ARGUMENTO: \{'La propuesta actual tiene un esfuerzo total de 260+729+489+103+11+28+213+302 = 2135, lo que deja un margen de 770. Sin embargo, excluye C67 (esfuerzo 806, importancia 780) y C10 (esfuerzo 326, importancia 780), que son artículos de alto valor para mí. Dado que tengo espacio, propongo añadir C10 y C61 ya está incluido. Espera, C61 ya está. Propongo añadir C10 y quizás C55 o C73. C10 (326, 780) es muy eficiente. C73 (668, 508) es menos eficiente. C55 (564, 583) es moderado. La mejor adición es C10. También C67 es muy valioso pero su esfuerzo es alto. Si añado C10, el esfuerzo sería 2135+326=2461. Quedan 444. Podría añadir B36 (41, 461) que es muy eficiente. Propongo: A43, A67, B42, B19, C61, A18, C43, C08, C10, B36.'\}} \\
\texttt{PROPUESTA: \{'A43', 'A67', 'B42', 'B19', 'C61', 'A18', 'C43', 'C08', 'C10', 'B36'\}} \\
            }
        }
    }
    & & \\ \\

    \theutterance \stepcounter{utterance}  
    & & \multicolumn{4}{p{0.6\linewidth}}{
        \cellcolor[rgb]{0.9,0.9,0.9}{
            \makecell[{{p{\linewidth}}}]{
                \texttt{\tiny{[GM$\rangle$P2]}}
                \texttt{ARGUMENTO: \{'La propuesta actual tiene un esfuerzo total de 260+729+489+103+11+28+213+302 = 2135, lo que deja un margen de 770. Sin embargo, excluye C67 (esfuerzo 806, importancia 780) y C10 (esfuerzo 326, importancia 780), que son artículos de alto valor para mí. Dado que tengo espacio, propongo añadir C10 y C61 ya está incluido. Espera, C61 ya está. Propongo añadir C10 y quizás C55 o C73. C10 (326, 780) es muy eficiente. C73 (668, 508) es menos eficiente. C55 (564, 583) es moderado. La mejor adición es C10. También C67 es muy valioso pero su esfuerzo es alto. Si añado C10, el esfuerzo sería 2135+326=2461. Quedan 444. Podría añadir B36 (41, 461) que es muy eficiente. Propongo: A43, A67, B42, B19, C61, A18, C43, C08, C10, B36.'\}} \\
\texttt{PROPUESTA: \{'A43', 'A67', 'B42', 'B19', 'C61', 'A18', 'C43', 'C08', 'C10', 'B36'\}} \\
            }
        }
    }
    & \\ \\

    \theutterance \stepcounter{utterance}  
    & & & \multicolumn{4}{p{0.6\linewidth}}{
        \cellcolor[rgb]{0.9,0.9,0.9}{
            \makecell[{{p{\linewidth}}}]{
                \texttt{\tiny{[GM$\langle$P2]}}
                \texttt{ARGUMENTO: \{'La propuesta incluye C10 y B36, lo cual es interesante. Sin embargo, al evaluar la importancia total para mí, noto que C43 (822) y C08 (808) son extremadamente valiosos, pero C42 (138) y C67 (780) están fuera. C67 tiene una importancia de 780 para mí, lo cual es muy alto. La propuesta actual tiene un esfuerzo de 260+729+489+103+11+28+213+302+326+41 = 2492. Quedan 413 de margen. Podría añadir C67 (806) si elimino artículos de menor importancia relativa para mí, como C43 (822) o C08 (808)? No, C43 y C08 son muy altos. Pero C67 (780) es comparable. Veamos: Si quito C43 (213) y C08 (302), ahorro 515. Añado C67 (806). El esfuerzo neto aumenta 291. 2492+291 = 2783. Quedan 122. Puedo añadir B36 (41) si no estaba? B36 ya está. Puedo añadir A18? A18 ya está. Puedo añadir C61? C61 ya está. La importancia ganada con C67 es 780. La perdida con C43 y C08 es 822+808=1630. No es rentable. Mi objetivo es maximizar la suma. La propuesta actual tiene: A43(65)+A67(121)+B42(262)+B19(783)+C61(461)+A18(583)+C43(822)+C08(808)+C10(389)+B36(668) = 4962. Es una suma muy alta. Acepto esta propuesta porque maximiza mi importancia total dentro del límite.'\}} \\
\texttt{ACEPTAR: \{'A43', 'A67', 'B42', 'B19', 'C61', 'A18', 'C43', 'C08', 'C10', 'B36'\}} \\
            }
        }
    }
     \\ \\

    \theutterance \stepcounter{utterance}  
    & & & \multicolumn{2}{p{0.3\linewidth}}{
        \cellcolor[rgb]{0.9,0.9,0.9}{
            \makecell[{{p{\linewidth}}}]{
                \texttt{\tiny{[GM$|$GM]}}
                \texttt{game successful} \\
            }
        }
    }
    & & \\ \\

    \theutterance \stepcounter{utterance}  
    & & & \multicolumn{2}{p{0.3\linewidth}}{
        \cellcolor[rgb]{0.9,0.9,0.9}{
            \makecell[{{p{\linewidth}}}]{
                \texttt{\tiny{[GM$|$GM]}}
                \texttt{end game} \\
            }
        }
    }
    & & \\ \\

\end{supertabular}
}

\end{document}
